\documentclass[11pt,a4paper]{article}
\usepackage[margin=24mm]{geometry}
\usepackage[T1]{fontenc}
\usepackage{lmodern}
\usepackage{amsmath,amssymb,amsthm}
\usepackage{graphicx,booktabs}
\usepackage{microtype}
\usepackage[hidelinks]{hyperref}
\hypersetup{pdftitle={A Geometric Characterisation of Prime Numbers Using Regular Polygons and Digons},pdfauthor={Pranay Manocha},pdfsubject={An expository geometric reformulation of elementary divisibility},pdfkeywords={prime numbers, regular polygons, digons, divisibility, divisor-counting function}}
\newtheorem{theorem}{Theorem}
\newtheorem{proposition}[theorem]{Proposition}
\newtheorem{corollary}[theorem]{Corollary}
\theoremstyle{definition}
\newtheorem{definition}[theorem]{Definition}
\setlength{\emergencystretch}{2em}
\raggedbottom
\title{A Geometric Characterisation of Prime Numbers\\Using Regular Polygons and Digons}
\author{Pranay Manocha\\[3pt]\small Independent researcher}
\date{5 October 2026\\[3pt]\small Expository note}

\begin{document}
\maketitle
\begin{abstract}
Fix a unit of length. We represent a positive integer $N$ by a regular
closed route with $k\ge2$ equal steps of integer length $s\ge1$, so that
$N=ks$. Ordinary regular convex polygons are supplemented by the
degenerate digon: a segment traversed out and back. For $N\ge2$, primality
is equivalent to every admissible route having unit step length. The
number of representation types is $R(N)=\tau(N)-1$, where $\tau$ counts
positive divisors. Thus $1$, primes and composites have respectively
zero, one and multiple types. This is an expository reformulation of
elementary divisibility, using at most two spatial dimensions.
\end{abstract}

\section{Motivation}
This construction arose from an attempt to develop a geometric language
for prime numbers as a preliminary step towards investigating additive
questions such as Goldbach's conjecture; no result concerning Goldbach's
conjecture is claimed here. Encoding factorisation geometrically does not
by itself resolve an additive question. The present purpose is to state
the correspondence precisely and treat the small integers uniformly.

\section{Definitions}
We work throughout with \emph{positive integers}, and fix a unit of length.
A prime is an integer greater than $1$ whose only positive divisors are
$1$ and itself. A composite is an integer greater than $1$ that is not
prime; $1$ is neither prime nor composite.

\begin{definition}
An \emph{admissible regular closed route} of total length $N$ has $k$ equal
straight steps of length $s$, where $k\ge2$ and $s\ge1$ are integers.
For $k\ge3$ it traverses the boundary of a regular convex $k$-gon once.
For $k=2$ it traverses a segment of length $s$ and retraces it to the
starting point; we call this a \emph{degenerate digon}. Its total
travelled length is $N=ks$.
\end{definition}

The digon is a convention for a closed route, not an ordinary
non-degenerate Euclidean polygon. Its two traversals coincide: the
geometric segment has length $s$, the travelled distance is $2s$, and
the final displacement is zero. No integer-coordinate condition is imposed.
A \emph{representation type} is determined solely by $(k,s)$.
Position, orientation, starting vertex and direction of traversal do not
create new types; in particular, rotations, translations and reflections
are not counted again. Star polygons and repeated circuits are excluded.

\newpage
\section{Geometric characterisation of primality}
Every pair of integers $k\ge2$, $s\ge1$ is realisable: for $k\ge3$,
inscribe a regular $k$-gon in a circle of radius
$s/(2\sin(\pi/k))$; for $k=2$, use the out-and-back route.
In particular, every $N\ge2$ has the unit-step type $(N,1)$.

\begin{theorem}\label{thm:prime}
For every integer $N\ge2$, $N$ is prime if and only if every admissible
regular closed route of total length $N$ has unit step length.
Equivalently, no such route has integer step length $s>1$.
\end{theorem}
\begin{proof}
If an admissible route has $s>1$, then $N=ks$ with $k,s\ge2$,
so $N$ is composite. Conversely, if $N$ is composite, write $N=ab$
with $a,b\ge2$. Choose $k=a$ and $s=b$. When $a=2$, use the digon;
when $a\ge3$, use a regular convex $a$-gon of side length $b$.
This is an admissible route with non-unit integer step length.
These two implications prove the equivalence. If $N$ is prime,
$k\mid N$ and $k\ge2$ force $k=N$, and hence $s=1$; its unique
representation type is $(N,1)$.
\end{proof}

The assumption $N\ge2$ is essential: $1$ admits no route, so the
universal unit-step condition would hold vacuously at $1$.

\section{The exceptional role of the digon}
If only ordinary polygons ($k\ge3$) were allowed, the same criterion
would work for $N\ge5$. Indeed, a composite $N\ge5$ has a factorisation
$N=ab$ with $a\ge b\ge2$ and $a\ge3$: otherwise $a=b=2$ and $N=4$.
At $4$, however, the only ordinary-polygon type is the unit square
$(4,1)$; the missing non-unit witness is $(2,2)$. Admitting the digon
provides this witness and also gives the prime $2$ its unit route $(2,1)$.
Thus $2=2\cdot1$ and $4=2\cdot2$ fit the same convention.

\paragraph{Equal turns.}
Each route can be followed using the same oriented exterior turn
$360^\circ/k$ after each edge, including at the closing vertex.
For $k=2$ the $180^\circ$ turn reverses direction along the segment;
for $k=3$ the $120^\circ$ turn produces an equilateral triangle.
This unifies the construction geometrically without enlarging the
admissible family to other turning angles or star polygons.

\section{Representation count and classification}
For a positive integer $N$, define
\[
 R(N)=\#\{(k,s)\in\mathbb Z^2:k\ge2,\ s\ge1,\ ks=N\}.
\]
Let $\tau(N)$ be the number of positive divisors of $N$ (also denoted
$d(N)$ in standard references~\cite{apostolDLMF27}).
\begin{proposition}\label{prop:count}
For every positive integer $N$, $R(N)=\tau(N)-1$.
\end{proposition}
\begin{proof}
Every divisor $k\mid N$ with $k\ge2$ gives exactly one type
$(k,N/k)$, and every type arises in this way. Of the $\tau(N)$
positive divisors, only $k=1$ is excluded.
\end{proof}
\begin{corollary}\label{cor:classification}
For positive integers $N$,
\[
 N=1\iff R(N)=0,\qquad
 N\text{ prime}\iff R(N)=1,\qquad
 N\text{ composite}\iff R(N)\ge2.
\]
\end{corollary}
\begin{proof}
The integer $1$ has no type. Every $N\ge2$ has exactly one unit-step
type, $(N,1)$; Theorem~\ref{thm:prime} shows that an additional type
exists exactly when $N$ is composite.
\end{proof}

\newpage
\section{Examples}
Table~\ref{tab:examples} lists all types for $1\le N\le12$.
The two factors have different roles, so, for example, $(3,4)$ and
$(4,3)$ are distinct types. Figure~\ref{fig:routes} makes the distinction
visible for $2\le N\le6$.

\begin{table}[!ht]
\centering
\caption{All representation types for the first twelve positive integers.}
\label{tab:examples}
\medskip
\renewcommand{\arraystretch}{1.15}
\begin{tabular}{@{}rll@{}}
\toprule
$N$ & Representation types $(k,s)$ & Classification\\
\midrule
% BEGIN REPRESENTATION TABLE
1 & none & unit\\
2 & $(2,1)$ & prime\\
3 & $(3,1)$ & prime\\
4 & $(2,2),(4,1)$ & composite\\
5 & $(5,1)$ & prime\\
6 & $(2,3),(3,2),(6,1)$ & composite\\
7 & $(7,1)$ & prime\\
8 & $(2,4),(4,2),(8,1)$ & composite\\
9 & $(3,3),(9,1)$ & composite\\
10 & $(2,5),(5,2),(10,1)$ & composite\\
11 & $(11,1)$ & prime\\
12 & $(2,6),(3,4),(4,3),(6,2),(12,1)$ & composite\\
% END REPRESENTATION TABLE
\bottomrule
\end{tabular}
\end{table}

\begin{figure}[!ht]
\centering
\includegraphics[width=\linewidth]{paper/figures/overview.pdf}
\caption{All types for $N=2,3,4,5,6$, at a common length scale.
Blue routes have $s=1$; amber routes have $s>1$, also indicated by
the labels. The digon's opposing arrows lie on a single segment,
which is traversed twice. Primes have one type; composites have more.}
\label{fig:routes}
\end{figure}

\newpage
\section{Dimensional scope}
The digon lies in one dimension, while ordinary regular polygons lie
in a plane. The construction therefore requires at most two spatial
dimensions. Every planar route embeds unchanged in $\mathbb R^n$ for
$n\ge2$, so the result imposes no upper bound on mathematical or
physical dimensions. Higher-dimensional extensions are left for
separate investigation; this note establishes only the planar case.

\section{Relation to existing geometric descriptions of primality}
Geometric descriptions of factorisation and primality already exist.
Rectangular arrays represent products by rows and columns; Way's
educational treatment explicitly connects their shapes with factors
and prime numbers~\cite{way2011arrays}. Polygonal and star-polygon
constructions also express divisibility and coprimality: Erfle describes
how a fixed jump among cyclically arranged vertices relates to whether
a star is traced in one circuit~\cite{erfleCoprime}.
Those constructions provide context; star polygons are not routes
admitted by our definition.

This note records an independently developed geometric reformulation
of elementary divisibility. The correspondence between geometric
representation and factorisation is not claimed to be a new result in
number theory. The purpose is to formulate the total-length construction
explicitly, remove the exception caused by ordinary polygon terminology
by admitting the degenerate digon, and derive a uniform representation
count. Historical priority for this exact formulation is not claimed;
whether this precise presentation has appeared before has not been
established. The references are contextual, not an exhaustive literature
or priority investigation.

\section{Conclusion}
Factorisation appears as alternative regular closed routes with equal
integer step lengths. Primes have one representation type, composites
have multiple types, and the total number is $\tau(N)-1$. The digon
makes the primality criterion uniform from $2$ upwards, while $1$ is
distinguished by having no representation. The result is a geometric
description of the prime/composite distinction, with its arithmetic
content made explicit.

\section*{Source and licence}
The source, vector diagrams and finite arithmetic checks are available
at \href{https://github.com/pranman/prime-geometric-characterisation}{the accompanying repository}.
The checks guard against transcription errors; the proofs above establish
the general results. Copyright \textcopyright\ 2026 Pranay Manocha.
The note and artwork are licensed under
\href{https://creativecommons.org/licenses/by/4.0/}{Creative Commons Attribution 4.0 International}.

\begingroup
\small
\raggedright
\bibliographystyle{plain}
\bibliography{references}
\endgroup
\end{document}
